\section{A half power for each new window}
\label{p13:half:section}

Relations with long selections require a different count.
The useful alternative is concrete: either several large odd
prime kernels determine another centre by congruences, or many
selected values have large smooth parts. Those smooth parts
then determine a rare centre by a modulus larger than \(N\).
Iterating this alternative reaches the square-root scale.

In this section the number of blocks and their widths are
fixed before \(N\). A block consists of distinct integers
in a translate of an interval of diameter \(O_D(B)\);
its values are comparable to \(N\), and its centre varies
in an interval of length \(O_D(N)\).
Repeated values in a normalized block are merged.
These conventions include original windows and the
fixed-width blocks obtained by compressing a rational pair.
Let \(S_z\) denote the image of its even selections.

\begin{theorem}[Conditional full-support count]
\label{p13:half:conditional}
Fix \(r\ge1\), \(s\ge0\), \(D>0\), and \(\varepsilon>0\).
Let \(V_{\rm old}\) be the space generated by all values
in \(s\) fixed source blocks. Uniformly in those sources,
the number of lists of distinct new centres \(z_1,\ldots,z_r\)
such that
\begin{equation}
 u_i\in S_{z_i}\setminus\{0\}\quad(1\le i\le r),
 \qquad \sum_i u_i\in V_{\rm old}
 \label{p13:half:condition}
\end{equation}
for some \(u_i\), is \(O_{r,s,D,C,\varepsilon}(N^{r/2+\varepsilon})\).
The \(u_i\) need not remain nonzero modulo \(V_{\rm old}\).
The source blocks may overlap or be dependent.
\end{theorem}

The condition asks for a nonzero projection in each new
effective space. It imposes neither affine compatibility
nor a separation condition on physical values.

For $H=C_DB$ larger than every block diameter, write
\[
 K_H(n)=\prod_{\substack{\ell>H\\v_\ell(n)\ {\rm odd}}}\ell.
\]
This is the odd prime kernel after the small primes have been removed.


\paragraph{Why the half power appears.}
For a fixed witness, short selections are counted by
Corollary~\ref{p13:short:old-space}.  A long selection either has many
large odd kernels, which determine another centre by congruences, or
has many large smooth parts, whose least common multiple exceeds the
length of the centre interval.  In the second case the centre lies in
a sparse deterministic set.

The complete proof in Section~\ref{p14:comp:half:iteration} iterates
these two alternatives.  At a stage indexed by $x_j\le1/2$, a surviving
source has at most $N^{1-x_{j-1}+o(1)}$ possibilities and its receiver
degree is at most $N^{x_j+o(1)}$.  Removing this source--receiver pair
and applying induction to the other $r-2$ blocks costs
$N^{r/2+x_j-x_{j-1}+o(1)}$.  Small fixed increments and a final
$N^{1/2+o(1)}$ exceptional source set give the theorem.
The induction is simultaneous in every fixed number of old sources,
so adding exposed blocks to that space introduces no dependence
assumption.

\begin{corollary}[Connected and primitive families at fixed order]
\label{p13:half:connected}
For every fixed \(r\ge2\), connected families of \(r\)
original windows have \(O(N^{r/2+\varepsilon})\) hosts.
Their independent mass, and in particular the independent
mass of primitive families, is \(O(N^{-r/2+\varepsilon})\).
\end{corollary}

\begin{proof}
A primitive family has a relation with nonzero projection
in each block, so Theorem~\ref{p13:half:conditional} applies
with no sources.
A connected family contains a minimal dependent subfamily;
expose it first. If \(A\) is the set already exposed,
connectedness implies
\[
 \left(\sum_{i\in A}S_i\right)
 \cap\left(\sum_{j\notin A}S_j\right)\ne\{0\}.
\]
Choose a minimal nonempty subset \(U\) of the complement
whose sum meets the old sum nontrivially.
A nonzero vector in that intersection has a representation
with all projections on \(U\) nonzero, by minimality.
The theorem counts these new blocks by
\(N^{|U|/2+\varepsilon'}\). Continue until the whole
family is exposed.
The group sizes sum to \(r\), and the number of possible
ordered partitions of the index set depends only on \(r\).
Their product count is \(N^{r/2+O_r(\varepsilon')}\).
Multiplying by \(\prod p_x\ll_C N^{-r}\) gives the mass.
\end{proof}

\begin{corollary}[Dependent pairs]
\label{p13:half:pairs}
The number of pairs of blocks with nonzero intersection
of their effective even-selection spaces is
\(O(N^{1+\varepsilon})\), uniformly for fixed logarithmic
widths and fixed normalizations.
\end{corollary}

\begin{proof}
Use Theorem~\ref{p13:half:conditional} with \(r=2,s=0\).
The two equal nonzero classes give the required relation.
\end{proof}

\paragraph{Sharpness and range.}
For \(r=2k\), take \(k\) pairs of adjacent windows,
with the pairs far apart. In each pair select the same
two consecutive values in the overlap. Their square
class is nonzero, since \(n(n+1)\) lies strictly between
\(n^2\) and \((n+1)^2\), and the equal projections cancel.
There are \(\gg N^k\) such lists.
Thus \(r/2\) is optimal for the general full-support
count at even orders. For \(k>1\) these examples can be
decomposable; they make no optimality claim for connected
or primitive subclasses. Every order and width in this
section remains fixed.
